\chapter{Cataract Symptoms}

\section*{Definition}
A cataract is an opacification of the crystalline lens that causes progressive, painless deterioration of visual quality; it is typically age-related but may be congenital, traumatic, or secondary to systemic disease. Symptoms correlate with the location and density of lens opacity.

\section*{What Happens in Your Body}
Accumulation of damaged lens proteins (crystallins) due to oxidative stress, ultraviolet radiation exposure, and metabolic derangements leads to protein aggregation and loss of lens transparency. This disrupts the passage of light, causing scatter, absorption, and irregular refraction.

\section*{Causes}
\begin{itemize}
 \item \textbf{Age-related:} Nuclear sclerotic, cortical, or posterior subcapsular subtypes developing after age 50
 \item \textbf{Traumatic:} Blunt or penetrating injury causing lens fiber disruption
 \item \textbf{Congenital:} Rubella, galactosemia, or genetic mutations (e.g., connexin genes)
 \item \textbf{Metabolic:} Diabetes mellitus accelerating sorbitol pathway accumulation
 \item \textbf{Drug-induced:} Chronic corticosteroid use (systemic, topical, or inhaled)
 \item \textbf{Radiation:} Ionizing radiation or prolonged ultraviolet B exposure
\end{itemize}

\section*{When to See a Doctor}

\begin{itemize}
 \item Progressive visual blurring interfering with reading, driving, or daily activities
 \item Glare and halos around lights, especially at night
 \item Frequent prescription changes for glasses
 \item Monocular diplopia or polyopia
 \item Decreased contrast sensitivity and color desaturation
\end{itemize}

\section*{Related Symptoms}
\begin{itemize}
 \item Decreased visual acuity (distance and near)
 \item Glare disability and photosensitivity
 \item Second sight (temporary myopic shift in nuclear sclerosis)
 \item Reduced color discrimination (blue-yellow defect)
 \item Strabismus or diplopia from anisometropia after unilateral progression
\end{itemize}
\textbf{Important:} Posterior subcapsular cataracts disproportionately impair near vision and cause severe glare in bright light despite relatively preserved distance acuity.

\section*{Self-Care / Home Management}
\begin{itemize}
 \item Improved ambient lighting and anti-glare sunglasses for outdoor activities
 \item Magnification aids for reading and near tasks
 \item Updating spectacle prescription as cataract progresses
 \item Pupil dilation with brighter lighting if posterior subcapsular cataract
 \item Avoiding night driving when glare becomes problematic
\end{itemize}

\section*{How Your Doctor Will Evaluate You}
\begin{itemize}
 \item Snellen visual acuity measurement with and without glare source
 \item Slit-lamp biomicroscopy after pupillary dilation (LOCS III grading system)
 \item Contrast sensitivity testing (Pelli-Robson chart)
 \item Potential acuity meter or laser interferometry to estimate post-surgical visual potential
 \item Specular microscopy for endothelial cell count if considering surgery
\end{itemize}

\section*{Treatment Options}
\begin{itemize}
 \item Observation for early cataracts with refractive compensation
 \item Cataract extraction (phacoemulsification) with intraocular lens implantation when visual function impairs quality of life
 \item Femtosecond laser-assisted cataract surgery as a precision alternative to manual phacoemulsification
 \item Management of underlying risk factors (glycemic control in diabetes, corticosteroid minimization)
 \item Nd:YAG laser capsulotomy for posterior capsule opacification developing months to years after surgery
\end{itemize}

\section*{References}
\begin{thebibliography}{99}
\bibitem{ref1} Liu YC, Wilkins M, Kim T, et al. "Cataracts." Lancet. 2017;390(10094):600--612.
\bibitem{ref2} Asbell PA, Dualan I, Mindel J, et al. "Age-Related Cataract." Lancet. 2005;365(9459):599--609.
\bibitem{ref3} Flaxman SR, Bourne RRA, Resnikoff S, et al. "Global Causes of Blindness and Distance Vision Impairment 1990--2020: A Systematic Review and Meta-Analysis." Lancet Glob Health. 2017;5(12):e1221--e1234.
\bibitem{ref4} Kohnen T, Baumeister M, Kook D, et al. "Cataract Surgery: Indications, Techniques, and Intraocular Lens Selection." Dtsch Arztebl Int. 2009;106(43):695--702.
\bibitem{ref5} Nizami AA, Gurnani B, Gul AC. "Cataract." In: StatPearls. Treasure Island (FL): StatPearls Publishing; 2024.
\end{thebibliography}
